%% kapa: exam=1 round=19 year=2008 period=1 form=A subject=civil_law subject_ko=민법 seq=26 number=26
%% source: Q-Net 감정평가사 기출문제 (cst00309) artlSeq=1001962 file=extracted/1008/1001962/2008년 제19회 감정평가사 1차시험 기출문제/1교시.hwp
%% conversion: kordoc 4.24.51 -> tools/gen1.py (auto); review=proofread-sonnet; flags=none
\begin{question}{26}
공동소유에 관한 설명으로 옳은 것은? (다툼이 있는 경우에는 판례에 의함)
\begin{choices}
  \item 부동산에 관하여 과반수의 지분을 가진 자는 다른 공유자와 협의가 없더라도 공유물의 관리에 관한 사항을 단독으로 결정할 수 있다.
  \item 공유물분할방법으로 공유물을 공유자 중 1인의 단독소유로 하고 그로 하여금 다른 공유자에게 그 지분의 가격을 배상시키는 방법은 허용되지 않는다.
  \item 공유물에 대한 소수의 지분권자는 그 공유물을 배타적으로 점유하는 다른 소수지분권자를 상대로 자신에게 공유물을 반환할 것을 청구할 수 없다.
  \item 합유자가 사망한 경우, 다른 특별한 약정이 없는 한 그 상속인이 합유지분을 상속한다.
  \item 법인 아닌 사단의 실질을 가진 교회가 사실상 분열된 경우 교회재산은 분열된 교회의 구성원들에게 각각 총유적으로 귀속된다.
\end{choices}
\end{question}
\answer{1}
